\documentclass[12pt, a4paper]{article}
\usepackage[utf8]{inputenc}
\usepackage[spanish]{babel}
\usepackage{amsmath, amssymb, amsthm}
\usepackage{geometry}
\geometry{margin=2.5cm}

\newtheorem{teo}{Teorema}
\newtheorem{ej}{Ejemplo}

\title{Determinantes y la Regla de Cramer}
\author{Álgebra Lineal}
\date{\today}

\begin{document}
	
	\maketitle
	
	\section{Introducción a los Determinantes}
	El determinante es una función que asigna un número escalar a una matriz cuadrada. Sea $A$ una matriz cuadrada de orden $n \times n$, su determinante se denota como $\det(A)$ o $|A|$.
	
	\subsection{Determinante de orden 2x2}
	Para una matriz de $2 \times 2$:
	\[
	A = \begin{pmatrix} a & b \\ c & d \end{pmatrix}
	\]
	Su determinante se calcula como:
	\[
	\det(A) = |A| = ad - bc
	\]
	
	\subsection{Determinante de orden 3x3 (Regla de Sarrus)}
	Para una matriz de $3 \times 3$, se pueden repetir las dos primeras columnas al final y sumar los productos de las diagonales principales restando las secundarias.
	
	\section{La Regla de Cramer}
	La regla de Cramer es un método matemático para resolver sistemas de ecuaciones lineales compatibles determinados (es decir, que tienen una solución única). 
	
	Utiliza los determinantes de la matriz de coeficientes del sistema y de las matrices que resultan de reemplazar una columna por los términos independientes.
	
	Dado un sistema de $n$ ecuaciones con $n$ incógnitas representado en forma matricial como $A \cdot X = B$, el valor de cada incógnita $x_i$ se calcula así:
	\[
	x_i = \frac{\det(A_i)}{\det(A)}
	\]
	Donde $A_i$ es la matriz resultante de sustituir la columna $i$-ésima de $A$ por el vector columna de los términos independientes $B$, y se exige que $\det(A) \neq 0$.
	
	\begin{ej}
		Resolver el siguiente sistema de ecuaciones lineales $2 \times 2$:
		\[
		\begin{cases}
			3x + 2y = 7 \\
			x - y = 1
		\end{cases}
		\]
		
		\textbf{Paso 1:} Identificar la matriz de coeficientes $A$ y el vector $B$:
		\[
		A = \begin{pmatrix} 3 & 2 \\ 1 & -1 \end{pmatrix}, \quad B = \begin{pmatrix} 7 \\ 1 \end{pmatrix}
		\]
		
		\textbf{Paso 2:} Calcular el determinante principal $\det(A)$:
		\[
		\det(A) = (3)(-1) - (2)(1) = -3 - 2 = -5
		\]
		
		\textbf{Paso 3:} Calcular $\det(A_x)$ sustituyendo la primera columna por $B$:
		\[
		A_x = \begin{pmatrix} 7 & 2 \\ 1 & -1 \end{pmatrix} \implies \det(A_x) = (7)(-1) - (2)(1) = -7 - 2 = -9
		\]
		
		\textbf{Paso 4:} Calcular $\det(A_y)$ sustituyendo la segunda columna por $B$:
		\[
		A_y = \begin{pmatrix} 3 & 7 \\ 1 & 1 \end{pmatrix} \implies \det(A_y) = (3)(1) - (7)(1) = 3 - 7 = -4
		\]
		
		\textbf{Paso 5:} Encontrar las soluciones para $x$ e $y$:
		\[
		x = \frac{\det(A_x)}{\det(A)} = \frac{-9}{-5} = \frac{9}{5}, \quad y = \frac{\det(A_y)}{\det(A)} = \frac{-4}{-5} = \frac{4}{5}
		\]
		\end{ej}
			
		\end{document}
